Q-learning bootstraps from the maximum of noisy value estimates, which is positively biased; double Q-learning removes this bias by decoupling action selection from evaluation. We reimplement tabular Q-learning, Double Q-learning, Weighted Double Q-learning and Maxmin Q-learning in numpy, vectorised over 1,000 independent runs per seed (5 seeds), and measure the estimated-minus-true value of the start state and the fraction of suboptimal actions on the Sutton--Barto maximisation-bias MDP and on random 20-state MDPs. On the trap MDP the textbook picture holds: the peak start-state overestimate is \rhval{main/q-learning/maxbias/bias_peak/mean} for Q-learning and \rhval{main/double-q-learning/maxbias/bias_peak/mean} for Double Q-learning, and the suboptimal-action fraction falls from \rhval{main/q-learning/maxbias/subopt_all/mean} to \rhval{main/double-q-learning/maxbias/subopt_all/mean}. On random MDPs the picture is mixed: every method \emph{underestimates} the start-state value at the end of training, Q-learning least (\rhval{main/q-learning/random20/bias_final/mean} versus \rhval{main/double-q-learning/random20/bias_final/mean} for Double Q-learning), and Double Q-learning takes more suboptimal actions (\rhval{main/double-q-learning/random20/subopt_all/mean} versus \rhval{main/q-learning/random20/subopt_all/mean}). A warm-start ablation, a step-size sweep and an exploration sweep are consistent with this being a finite-time effect of slower per-table learning and of adaptive data collection: with a warm start Double Q-learning takes fewer suboptimal actions, and with uniformly random behaviour the sign of Q-learning's bias is positive. These are small tabular experiments with one MDP family and fixed hyper-parameters; they do not transfer directly to deep RL.
